% !TeX root = ../collection-solutions.tex

\titledquestion{Social comparisons} In the model of \citet{drechsel2025falling-behind}, the optimal house of a given type~$i$
is given by $h^i = \kappa_2 y^i + \tilde \phi \tilde h^i$.

\begin{parts}

    \part Show that $\vec h = \kappa_2 (I + L) \vec y$ where $L$ is the social externality matrix.

    \begin{solution}
        \begin{itemize}
            \item Stack the equations: $\vec h = \kappa_2 \vec y + \tilde \phi \tilde{\vec h}$
            \item Use the definition $\tilde{\vec h} = G \vec h$: 
            $$\vec h = \kappa_2 \vec y + \tilde \phi G \vec h$$
            \item solve
            $$(I - \tilde \phi G)\vec h = \kappa_2 \vec y \implies \vec h = \kappa_2 (I - \tilde \phi G)^{-1} \vec y $$
            \item use that the Leontief inverse has a sum representation: $$(I - \tilde \phi G)^{-1} = \sum_{k=\textcolor{red}{0}}^\infty (\tilde \phi G)^k = I + \underbrace{\sum_{k=\textcolor{red}{1}}^\infty (\tilde \phi G)^k}_{=: L}$$
            \item Plug in: $\vec h = \kappa_2 (I + L) \vec y$
        \end{itemize}
    \end{solution}

    \part Compute $\frac{\partial}{\partial y^j} h^i$ (where $j \neq i$).

    \begin{solution}
        From (a) we know that $h_i = \kappa_2 (y_i + \sum_{\textcolor{blue}{k} = 1}^n  
        L_{i\textcolor{blue}{k}}
        y_{\textcolor{blue}{k}})
        $, where $L_{ik}$ is the $(i, k)$-entry of the social externality matrix $L$.

        Thus, $\frac{\partial}{\partial y^j} h^i = \kappa_2 L_{ij}$.
    \end{solution}

    \part Interpret this result and discuss its relevance.

    \begin{solution}
        \begin{itemize}
            \item   $L_{ij}$ is the externality that $j$ exerts on $i$. It is positive, if $\phi > 0$ and if $i$ cares about $j$'s house (directly or indirectly).

            \item If $j$'s income goes up, and $j$ exerts an externality on $i$, then $i$ will spend more on housing (and less on consumption).
    
            \item Such behavior is well-documented empirically. It can help explain why household debt has increased so much between 1980 and 2007.  
        \end{itemize}
      \end{solution}

\end{parts}